\documentclass[11pt]{article}
\usepackage[a4paper,margin=25mm]{geometry}
\usepackage{amsmath,amssymb,amsthm,booktabs,xurl,hyperref}
\hypersetup{hidelinks}
\setlength{\emergencystretch}{3em}
\newtheorem{theorem}{Theorem}
\newtheorem{lemma}[theorem]{Lemma}
\newtheorem{proposition}[theorem]{Proposition}
\newcommand{\Q}{\mathbb Q}
\newcommand{\Z}{\mathbb Z}
\DeclareMathOperator{\Tr}{Tr}
\DeclareMathOperator{\Nm}{N}
\DeclareMathOperator{\sqf}{sqf}
\title{Congruent-number poles in the complete combined quartic systems}
\author{Z diophantine in Q project}
\date{24 November 2026}
\begin{document}
\maketitle

\begin{abstract}
We replace all 39 or 44 private non-CM pole blocks in the complete
combined general and specified systems by congruent-number poles.
The resulting systems retain 963 variables and 818 quadratic rows,
or 1140 variables and 991 rows. Exactly 117 or 132 rows change.
Their integer-input interpretation uses Bhargava--Burungale's
Theorem~2.3 instead of the former non-CM pole realization input.
We prove the simultaneous local preparation with the actual shared
reference point, scalar point, and global contact gauge, and give
all six physical coordinate maps. The other arithmetic interfaces
and the external integer-coding theorem remain premises. This is
a reduction of assumptions, not a further reduction of variables
or a bijection between the old and new rational zero sets.
\end{abstract}

\section{Statement and precise scope}

The predecessors are the complete contact-denominator systems in
\cite{v103}. They contain thirteen integer predicates in the general
route and eleven in the specified route. The general route uses two
comparison poles and one scalar pole per predicate; the specified
route uses three comparison poles and one scalar pole. Each pole
has private auxiliary coordinates, even though the physical point
data and contact gauge have already been shared.

We retain the predecessor's elliptic carrier over $K=\Q(\sqrt2)$,
point-chart interpretation, formal-logarithm estimates, height bounds,
finite-family prime supplier, all-sign reduction-order coprimality,
and compatible finite-place and compact-real recurrence. The
general defining parameter retains its ineffective sufficiently
large threshold; the specified parameter is unchanged. The later
published integer-coding theorem recorded in the full tables is an
external premise, and is not independently proved here. Historical
papers in the dependency chain sometimes discuss an older emended
source: those historical source obligations are not substituted for
the current table's source-coding metadata.

The one arithmetic premise removed here is realization of the old
non-CM pole predicate. In particular, the new construction does not
invoke the old Theorem~5.3, Route~K, its certified-rank input, or the
project's Proposition~N. It invokes the rational-point conclusion of
Bhargava--Burungale~\cite[Theorem~2.3]{BB}. We verify the translation
and every additional preparation condition below. Finite coefficient
checks establish neither that theorem nor the other retained
arithmetic interfaces.

\begin{theorem}\label{main}
Under the retained non-pole arithmetic interfaces and the cited
Bhargava--Burungale result, the new systems have the same intended
integer-input projection as the predecessor systems. They have,
respectively, $(N,r)=(963,818)$ and $(1140,991)$, where $N$ counts
rational variables and $r$ counts quadratic residuals. Their sums
of squared residuals have degree exactly four. No complete old/new
rational zero-set bijection is asserted.
\end{theorem}

\section{The elementary obstruction and the external point input}

For a nonzero rational $d$, consider
\begin{equation}\label{cm}
 e^2+dp^2=1,\qquad uv=e(1-d),\qquad u^2-v^2=dq^2.
\end{equation}
The letters $p,q$ in this display are rational auxiliaries, not primes.

\begin{lemma}\label{obstruction}
Every rational solution of \eqref{cm} has square reduced denominator
for $d$: each negative valuation of $d$ is even.
\end{lemma}
\begin{proof}
Suppose $v_\ell(d)=\nu<0$ is odd. In the first equation the two
summand valuations have different parity whenever $p\ne0$. Neither
can be negative, and $e$ is a unit. If $p=0$, then $e=\pm1$, with
the same conclusion. The middle equation gives $uv\ne0$ and
$v_\ell(u)+v_\ell(v)=\nu$. If $q=0$, the last equation implies
$v_\ell(u)=v_\ell(v)$. If $q\ne0$ and these two valuations differ,
the left side has even valuation whereas the right side has odd
valuation. Thus they again coincide. Their sum is even, a
contradiction. This includes $\ell=2$.
\end{proof}

\begin{proposition}[Prepared existence]\label{existence}
Assume the rational-point assertion of \cite[Theorem~2.3]{BB}.
If $0<d<1$, $d$ has square reduced denominator,
$v_2(d)\ge4$, and $d\in7\Q_2^{\times2}$, then \eqref{cm} has
rational solutions.
\end{proposition}
\begin{proof}
Write $d=A/B^2$ in lowest terms with $A,B>0$. The positive
squarefree representatives $j=\sqf(d)$ and $k=\sqf(1-d)$ are
coprime, because $A$ and $B^2-A$ are coprime. They are odd,
$j\equiv7\pmod8$, and $k\equiv1\pmod8$. At a prime dividing
$j$, $1-d$ is a unit congruent to one, so $k$ is a local square.
At a prime dividing $k$, $d$ is a unit congruent to one, so $j$
is a local square. These verify the mutual-square conditions in
the cited theorem. If $k=1$, put $1-d=r^2$ and take
$e=1$, $p=q=0$, $u=v=r$; this allowed zero branch already proves
the assertion.

For $k>1$, the theorem supplies a positive $n$ and points of
nonzero ordinate on
\[
 E_n:y^2=x^3-n^2x,\qquad E_{nk}:y^2=x^3-n^2k^2x,
\]
whose abscissas lie in $j\Q^{\times2}=d\Q^{\times2}$.
Writing $1-d=kc^2$, scale the second point by
$(x,y)\mapsto(c^2x,c^3y)$. Thus there are points
\[
 (da^2,aU)\in E_n(\Q),\qquad
 (db^2,bV)\in E_{n(1-d)}(\Q),\qquad abUV\ne0.
\]
Their equations are
\[
 U^2=d^3a^4-dn^2,\qquad
 V^2=d^3b^4-dn^2(1-d)^2.
\]
Define
\begin{equation}\label{pointrecipe}
 e=\frac{n}{da^2},\quad p=\frac{U}{d^2a^2},\quad
 u=\frac ba,\quad v=\frac{e(1-d)}u,\quad
 q=\frac{V}{d^2ab}.
\end{equation}
Direct substitution gives all three equations \eqref{cm}.
The independent verifier checks these identities as rational
identities in the two curve residuals. No density assertion or
additional nonzero condition on the pole auxiliaries is needed:
the complete tables retain private affine pole witnesses, with no
finite-algebra encoder or guard dividing by $p$, $q$, or $e$.
\end{proof}

Put $d=7b/9$. Clearing the constant denominator gives the new
physical pole equations
\begin{equation}\label{newphysical}
 9e^2-9+7bp^2=0,\quad
 9e-7eb-9uv=0,\quad
 9u^2-9v^2-7bq^2=0.
\end{equation}
By Lemma~\ref{obstruction}, these exclude every negative odd
valuation of $b$ outside $\{7\}$: the change at $3$ is by $-2$,
which preserves oddness and negativity. The prime $7$ is already
excluded from the consumed height radicals. For example, the
retained fixed height set in \cite[Section~5]{v62} is
\[
 S=\{2,3,5,7,11,13,17,37,41,2399\},
\]
and \cite[Section~5]{v82} also explicitly includes $7$ among its
existing exclusions. No contact-height constant changes.

For completeness, the old preparation gives square-denominator
inputs $b>0$, which are $2$-adic squares of valuation at least four
and $3$-adic squares of valuation at least two. We now arrange
simultaneously that $b<1$ and $v_7(b)\ge0$. Then $0<d<7/9$,
$v_2(d)\ge4$, and $d\in7\Q_2^{\times2}$. At $3$,
$v_3(d)=v_3(b)-2\ge0$; at $7$, $v_7(d)=v_7(b)+1\ge1$.
All other negative valuations remain even. Hence
Proposition~\ref{existence} applies to every actual input.

\section{Simultaneous preparation with shared physical points}

Fix the finite tuple of integer inputs. The retained normalizers
$R_i>100$ are fixed before the common point choices. Their positivity
also holds on every rational input: in the source notation
\[
 \beta=51X^2+2504,\qquad
 r=\beta(\beta-1)(\beta-2401)(32\beta+343)>0,
\]
and $R=149C^2r/(1+149^2r)$ or $R=A_{\rm a}r/(1+25r)$.
The lower bound $R_i>100$ on integer witnesses is the one already
used in \cite{v82}. We keep every scalar pole private, while
using one reference point $H$, one scalar point $Q_2$, and the
common Lucas roles furnished by the finite-family construction
\cite{common,global}. The contact-progression and endpoint changes
\cite{progression,aggregate,v102,v103} retain these physical choices.

\subsection{One scalar depth for every input}

Write $H=[h]P_0$, $Q_2=[a_0]P_0$, $a=2a_0$, and
$\ell=a/h$. The letter $\ell$ here is a rational index ratio.
Let $z=x/y$ be the fixed formal parameter, $f=1/x(H)$, and
$Z=2z(Q_2)/z(H)$. For each normalizer the actual scalar quantities
are exactly those of \cite[Sections~2 and 7]{v82}:
\begin{equation}\label{scalar}
 J_i=\frac{Z^2}{R_i^5}+\frac{f^2}{R_i^4},\quad
 \tau_i\Tr\left(\frac{Z^2}{R_i^5}\right)=\Nm(J_i),\quad
 b_{s,i}=K_i\tau_i^2,\quad K_i=R_i(R_i+9792).
\end{equation}
The finite preparation set is the union, over all inputs, of the
odd-valuation supports of $K_i$, the original exceptional preparation
primes, and $7$. This is a finite union after all $R_i$ are fixed.

At $7$, put $r_i=v_7(R_i)$ and $k_i=v_7(K_i)$. Choose one integer
\begin{equation}\label{tchoice}
 t\ge\max\left(1,\max_i\left\lceil\frac{10r_i-k_i}{4}\right\rceil\right),
\end{equation}
increased if an original preparation requirement asks for a larger
depth. Choose one common reference point deep enough that, at both
places $w$ above $7$, with $d_w=v_w(z(H))$ and some fixed $q\ge2$,
\begin{equation}\label{deep}
 4d_w-c_7\ge q,\qquad r_i+4d_w-2t\ge q\quad\hbox{for every }i.
\end{equation}
Here $c_7$ is a fixed formal-expansion denominator loss. The curve
$y^2=x^3-(11-6\sqrt2)x$ need not have good reduction at each place
above $7$. On a sufficiently deep formal subgroup its convergent
short-Weierstrass logarithm has leading term $z$ and next possible
term of order five. This gives, with a fixed loss on that subgroup,
\[
 f=z(H)^2\cdot\text{unit},\qquad
 Z=\ell(1+\epsilon_w),\qquad v_w(\epsilon_w)\ge4d_w-c_7.
\]
This is the same fixed-prime formal preparation used in
\cite[Section~7]{v82}; no good-reduction assertion at $7$ is needed.

Choose the common scalar index with the exact open condition
$v_7(a/h)=t$. Chinese remaindering combines this with all original
prepared-place conditions, $a\in6\Z$, and the fixed component and
reduction orders. The retained recurrence supplies the simultaneous
real return. Conditions \eqref{deep} imply that $J_i$ is
$\ell^2R_i^{-5}$ times a field unit congruent to one, and its trace
denominator in \eqref{scalar} is $2\ell^2R_i^{-5}$ times a rational
unit. Consequently
\begin{equation}\label{depth}
 v_7(\tau_i)=2t-5r_i,\qquad
 v_7(b_{s,i})=k_i+4t-10r_i\ge0
\end{equation}
for every input, using the same $H$ and $a_0$.

At any other prepared rational prime $p$, the identical calculation
from the source gives
\[
 v_p(\tau_i)=2t_p-5v_p(R_i)-v_p(2).
\]
Choose the one common lower requirement for $t_p$ as the maximum
over the finite input list, then increase the common reference
depth to allow it. Thus the previous squareclasses and depths at
$2,3,17$, and all odd-support integrality conditions, hold
simultaneously. This does not select a separate reference or scalar
point for different inputs.

For clarity, no new denominator prime of $\tau_i$ escapes the
argument. If $R_i=u/v$ is reduced, then
\[
 K_i=\frac{u(u+9792v)}{v^2},\qquad
 \gcd\bigl(u(u+9792v),v\bigr)=1.
\]
Every odd valuation of $K_i$ is positive and lies in the preparation
set. At such a prime we require $\tau_i$ integral. Outside that set
$v_p(K_i)$ is even, so $v_p(K_i\tau_i^2)$ is even regardless of
any denominator of $\tau_i$. This proves square denominators for
all actual scalar inputs.

The single real box
\[
 |f^{(j)}|\le1,\qquad 1\le|Z^{(j)}|\le3\quad(j=1,2)
\]
works for every $R_i$. The existing estimate gives
\[
 0<b_{s,i}\le\frac{24482500}{R_i^4}<1.
\]
After $a_0$ is fixed, shrink the common reference progression to
preserve these real and finite open conditions. This happens before
the comparison primes are chosen. Later high-precision reference
and common-numerator choices remain inside this progression.

\subsection{The actual common contact gauge at the added prime}

After the private root differences, common scalar and Lucas roles,
and aggregate shifts are fixed, there are $m=26$ comparison
contacts in the general system or $m=33$ in the specified system.
The progression includes the retained sixteenth endpoint. Write
its common contact parameter as $\lambda_c$ and all already fixed
shifts as $\delta_{ij}$. The shared gauge is the actual sum
\[
 \rho=\sum_{i,j}q_{ij}(\lambda_c+\delta_{ij}).
\]
At each prepared finite place, including $7$, and at real infinity,
the reflected contact functions satisfy the existing asymptotic
\[
 q_{ij}=\kappa\lambda_c^{-2}+O(\lambda_c^{-3}),\qquad
 \rho=m\kappa\lambda_c^{-2}+O(\lambda_c^{-3}),\qquad\kappa\ne0.
\]
In particular,
\begin{equation}\label{gaugeasymptotic}
 \frac{(\rho-q_{ij})^2}{q_{ij}}
       =(m-1)^2\kappa\lambda_c^{-2}+O(\lambda_c^{-3})\longrightarrow0.
\end{equation}
No division by $m$ or $m-1$ is used. All physical tests are bounded
on the chosen common point neighborhoods, and each nonconstant
detector term contains a positive power of $\rho$. Thus
\[
 W_i\longrightarrow1,\qquad
 D_i=\frac{\Nm(W_i-1)}{\Tr(W_i)}\longrightarrow0,
\]
with trace tending to $2$, a $7$-adic unit. The physical gate and
pole input are
\begin{equation}\label{comparisoninput}
 g_{ij}=1+D_i^2+\frac{(\rho-q_{ij})^2}{q_{ij}},\qquad
 b_{ij}=\left(\frac{204}{409}\right)^2\frac{g_{ij}^2}{1+q_{ij}}.
\end{equation}
Consequently every $b_{ij}$ tends to $(204/409)^2$, a $7$-adic
unit and a real number in $(0,1)$. A single sufficiently small
neighborhood of $1/\lambda_c=0$ at $7$, and a sufficiently large
real slope, give $v_7(b_{ij})=0$ and $0<b_{ij}<1$ for the whole
finite list. The previous preparations at $2,3,17$ persist.

This limiting argument uses the sum of all actual contact weights,
not the former two- or three-term sum of one input block. The
retained finite-family supplier permits fixed finite divisibility,
finite avoidance, and an arbitrarily large real slope; its
high-precision common-index construction preserves all prescribed
point neighborhoods. Apply it after adjoining $7$ to the fixed
preparation set and to the fresh-prime avoidance set. The source
separation argument, including the common spacing and endpoint,
is unchanged because all shifts were fixed before this step.

At a fresh prime of a contact $(i,j)$, its own weight has valuation
$-1$ and all other weights are integral. Hence $v_p(\rho)=-1$,
$\rho-q_{ij}$ is integral, and the original synchronized formal
depths give $D_i\equiv0$, $g_{ij}\equiv1$, and $v_p(b_{ij})=1$.
For a foreign pole, its contact factor has even valuation and the
gate changes it only by a rational square. Its valuation is
therefore even, even if the foreign gate has a pole. These are
exactly the global-gauge parity arguments of \cite{global}.
They prove square reduced denominators for every actual comparison
input while retaining the new integrality at $7$.

All shared physical choices are now fixed. Apply
Proposition~\ref{existence} independently to the 39 or 44 private
pole inputs, and fill their private stored coordinates as below.
No scalar, root, contact, or numerator choice is revisited after
these pole witnesses have been chosen.

\section{Six exact physical charts and the complete row change}

The old physical triple was
\begin{equation}\label{oldphysical}
 3e^2-34e-289+17bp^2=0,\quad
 17e-eb-17uv=0,\quad
 u^2+34uv-867v^2-17bq^2=0.
\end{equation}
For every pole let $E,B,U,V,P,Q$ denote its six stored coordinates.
There are nonzero scale quantities $s_1,s_2,S$ with $S^2=s_1s_2$.
The physical dictionary is
\begin{equation}\label{dictionary}
 E=s_1e,\quad B=s_2b,\quad U=Su,\quad V=Sv,\quad
 P=\frac{s_1^2}{s_2}p^2,\quad Q=s_1q^2.
\end{equation}
The two retained square rows are $p_*^2=\gamma_pP$ and
$q_*^2=\gamma_qQ$. All six charts are listed here; $R$ is the
private positive normalizer, $\epsilon=\rho^{-1}$ is the guarded
common inverse gauge, and $\sigma$ is the private scalar scale.

\begin{center}\small
\begin{tabular}{@{}lccccccc@{}}
\toprule
Chart & $s_1$ & $s_2$ & $S$ & $\gamma_p$ & $\gamma_q$
 & $p_*/p$ & $q_*/q$\\
\midrule
General 0 & $1$ & $\epsilon^2$ & $\epsilon$ & $R^2$ & $R^2$
 & $R/\epsilon$ & $R$\\
General 1 & $1$ & $\epsilon^2$ & $\epsilon$ & $1$ & $R^2$
 & $1/\epsilon$ & $R$\\
Specified 0 & $R^2$ & $1$ & $R$ & $\epsilon^2$ & $R^2$
 & $\epsilon R^2$ & $R^2$\\
Specified 1 & $R^2$ & $1$ & $R$ & $\epsilon^2$ & $\epsilon^2$
 & $\epsilon R^2$ & $\epsilon R$\\
Specified 2 & $\epsilon^2$ & $1$ & $\epsilon$ & $\epsilon^2$ & $\epsilon^2$
 & $\epsilon^3$ & $\epsilon^2$\\
Scalar & $\sigma$ & $\sigma$ & $\sigma$ & $\sigma$ & $\sigma$
 & $\sigma$ & $\sigma$\\
\bottomrule
\end{tabular}
\end{center}
The existing square registers for $R^2$ and $\epsilon^2$ are used
in the literal rows. The guard $\rho\epsilon=1$ gives
$\epsilon\ne0$. The scalar inverse row
$\sigma\widehat\tau=R^4$ gives $\sigma\ne0$, since $R>0$.
Thus every map and inverse in the table divides only by quantities
already forced nonzero; no pole auxiliary itself is inverted.

In these stored coordinates the old three rows are
\begin{align*}
 3E^2-34Es_1-289s_1^2+17BP&=0,\\
 17Es_2-EB-17UV&=0,\\
 U^2+34UV-867V^2-17BQ&=0.
\end{align*}
Replace them by
\begin{equation}\label{storednew}
\begin{aligned}
 9E^2-9s_1^2+7BP&=0,\\
 9Es_2-7EB-9UV&=0,\\
 9U^2-9V^2-7BQ&=0.
\end{aligned}
\end{equation}
The literal residuals have degree at most two, because $s_1,s_2$
are existing coordinates or constants. Substitution of
\eqref{dictionary} turns both old and new triples into their
physical triples multiplied by $s_1^2,S^2,S^2$, respectively.
The unchanged square rows recover rational $p$ and $q$ by the
last two columns of the table, including $p=0$ or $q=0$.

\subsection{Coordinate and occurrence audit}

Each input uses its original local coordinate prefix. In the
reciprocal input block the prefix is \texttt{int\_i\_reciprocal};
all others use \texttt{int\_}\emph{label}. The index dictionary is
\begin{center}\small
\begin{tabular}{@{}lrrrrrrrr@{}}
\toprule
Role & $E$ & $U$ & $V$ & $B$ & $P$ & $Q$ & $p_*$ & $q_*$\\
\midrule
Comparison 0 & 3&4&66&64&65&67&5&6\\
Comparison 1 & 7&8&75&73&74&76&9&10\\
Comparison 2 (specified) &166&167&176&174&175&177&168&169\\
Scalar &11&12&84&82&83&85&13&14\\
\bottomrule
\end{tabular}
\end{center}
The normalizer and its square are local indices $90,89$; scalar
$\sigma$ is index $81$. In the root block, the common
$(\rho,\epsilon,\epsilon^2)$ are $(159,160,161)$ in the general
system and $(162,163,164)$ in the specified system. These are
names inherited from earlier charts, not final positions after
renumbering. The verifier resolves them through the full table.

The independent checker imports no builder code. It reconstructs
each old and new pole row from this index dictionary and the six
physical maps, pins all 166 square rows, and verifies that every
other row is literally identical. It also checks both common
gauge guards, each scalar inverse and input equation, all variable
occurrences, coefficient rings, degrees, source metadata, and
file hashes. Its totals are
\begin{center}
\begin{tabular}{@{}lrrrrrr@{}}
\toprule
System & Variables & Rows & Poles & Changed & Unchanged & Square rows\\
\midrule
General &963&818&39&117&701&78\\
Specified &1140&991&44&132&859&88\\
\bottomrule
\end{tabular}
\end{center}
The full coefficient tables, not only the changed blocks, are
included. Their sum of squares vanishes exactly when all listed
quadratic residuals vanish over $\Q$. There is a nonzero quadratic
residual, so its degree is exactly four, since squared real leading
forms cannot cancel.

\section{Completing the implication and limits of the checkpoint}

For soundness, restore the unchanged contact, point, and normalizer
interpretations through the previous guarded charts. Each new
private pole has the physical interpretation \eqref{newphysical},
so Lemma~\ref{obstruction} supplies the same obstruction at every
consumed prime. The global-gauge valuation proof, scalar extraction,
and height comparison use only that obstruction, not the old
non-CM coefficients. They therefore give the same integrality
conclusions for the thirteen or eleven distinguished inputs.
If a contact lies at its allowed midpoint boundary, use its
retained cleared rows, as in \cite{global}; the quotient
\eqref{comparisoninput} is used for completeness and at consumed
primes, never to discard that boundary.

For completeness, Section~3 prepares all actual inputs with the
same common physical points and gauge. Proposition~\ref{existence}
then supplies private rational poles, and Section~4 fills their
stored coordinates. The remaining rows are those of the inherited
simultaneous construction. This proves Theorem~\ref{main} under
its stated premises. It neither maps arbitrary old pole witnesses
to new ones nor identifies their full zero sets.

The finite audit verifies the algebra and gives one prepared local
example:
\[
 b=\frac{144}{841},\quad d=\frac{112}{841},\quad
 (e,p,u,v,q)=\left(\frac34,\frac{29}{16},
                 \frac{27}{29},\frac{81}{116},\frac{27}{16}\right).
\]
This is a pole witness, not a witness for the complete integer
system and not computational evidence replacing the cited arithmetic
theorem. The checkpoint reduces arithmetic dependencies while
preserving the existing variable bounds. It makes no global
minimality claim and no claim that the quartic target $263$ has
been reached.

\begin{thebibliography}{9}
\bibitem{BB} M. Bhargava and A. Burungale,
\emph{A Diophantine definition of $\Z$ in $\Q$},
arXiv:2611.16130v2, Theorem~2.3 and Lemma~2.2.
\url{https://arxiv.org/html/2611.16130v2}.
\bibitem{v82} \emph{A scalar-point sharing construction in 163 and
179 variables}, v82, especially Sections~2, 5, and 7.
\nolinkurl{papers/quartic/scalar_point163_179_20261120_v82.tex}.
\bibitem{v62} The fixed three-contact height construction, v62,
Section~5, \emph{Arithmetic inputs and the finite exceptional set}.
\nolinkurl{papers/quartic/quartic229_20261117_v62.tex}.
\bibitem{common} \emph{Common physical points in a combined quadratic
system}, including the common-numerator all-sign CRT construction.
\nolinkurl{undecidability/papers/common_points_quartic_checkpoint.tex}.
\bibitem{global} \emph{One contact gauge in a combined quadratic
system}, Section~2.
\nolinkurl{undecidability/papers/global_gauge_quartic_checkpoint.tex}.
\bibitem{progression} \emph{A common arithmetic progression of contact
arguments and shared coefficients of their centered powers}.
\nolinkurl{undecidability/general_contact_v94/affine_contact_progressions.tex}.
\bibitem{aggregate} The aggregate-spacing checkpoint.
\nolinkurl{undecidability/aggregate_spacing_v100/aggregate_spacing.tex}.
\bibitem{v102} \emph{Two retained endpoint powers replace two contact
helpers}.
\nolinkurl{undecidability/contact_endpoint_v102/contact_endpoint_reuse.tex}.
\bibitem{v103} \emph{One polynomial denominator for a progression of
contact equations}.
\nolinkurl{undecidability/quadratic_reduction_v103/contact_denominator_compression.tex}.
\end{thebibliography}
\end{document}
